\documentclass[10pt]{article}
\usepackage[utf8]{inputenc}
\usepackage[a4paper,height=20cm,width=13cm]{geometry}
\usepackage{amssymb}
\usepackage{dsfont}
\usepackage{calc}
\usepackage{graphicx}
% \usepackage{pst-node}
\usepackage{fourier}
\usepackage{euscript}
\usepackage{amsmath,amsthm}
\def\phi{\varphi}
\def\theta{\vartheta}
\def\P{\EuScript P}
\def\M{\EuScript M}
\def\B{\EuScript B}
\def\C{\EuScript C}
\def\D{\EuScript D}
\def\U{\EuScript U}
\def\S{\EuScript S}
\def\sm{\smallsetminus}
\def\niff{\nleftrightarrow}
\def\ZZ{\mathds Z}
\def\NN{\mathds N}
\def\PP{\mathds P}
\def\QQ{\mathds Q}
\def\RR{\mathds R}
\def\<{\langle}
\def\>{\rangle}
\def\E{\exists}
\def\A{\forall}
\def\0{\varnothing}
\def\imp{\rightarrow}
\def\iff{\leftrightarrow}
\def\IMP{\Rightarrow}
\def\IFF{\Leftrightarrow}
\def\range{\textrm{im}}
\def\Mod{\textrm{Mod}}
\def\Aut{\textrm{Aut}}
\def\Th{\textrm{Th}}
\def\acl{\textrm{acl}}
\def\dom{{\rm dom}}
\def\eq{{\rm eq}}
\def\tp{\textrm{tp}}
\def\equivL{\stackrel{\smash{\scalebox{.5}{\rm L}}}{\equiv}}

\def\cnonfork{\mathbin{\raise1.8ex\rlap{\kern0.6ex\rule{0.6ex}{0.1ex}}\rlap{\kern1.1ex\rule{0.1ex}{1.9ex}}\raise-0.3ex\hbox{$\smile$} } }

\newcommand{\labella}[1]{{\sf\footnotesize #1}\hfill}
\renewenvironment{itemize}
  {\begin{list}{$\triangleright$}{%
   \setlength{\parskip}{0mm}
   \setlength{\topsep}{0mm}
   \setlength{\rightmargin}{0mm}
   \setlength{\listparindent}{0mm}
   \setlength{\itemindent}{0mm}
   \setlength{\labelwidth}{3ex}
   \setlength{\itemsep}{0mm}
   \setlength{\parsep}{0mm}
   \setlength{\partopsep}{0mm}
   \setlength{\labelsep}{1ex}
   \setlength{\leftmargin}{\labelwidth+\labelsep}
   \let\makelabel\labella}}{%
   \vspace*{-.3\baselineskip}
  \end{list}}
%\def\ssf#1{\textsf{\small #1}}
\newcounter{ex}
\newenvironment{exercise}{\bigskip\addtocounter{ex}{1}\textbf{Esercizio \theex.\ }}{}
\pagestyle{empty}
\parindent0ex
\parskip1ex
\raggedbottom
\def\nsR{{}^*\!\RR}
\def\ssf#1{\textsf{#1}}
\renewcommand{\baselinestretch}{1.3}


\usepackage{fancyhdr}
\pagestyle{fancy}
\lfoot{Teoria dei Modelli a.a.~2026/27}
\rhead{}
\rfoot{{\small\thepage}}
\cfoot{}



%
% INDICARE NOME E COGNOME DI TUTTI GLI AUTORI
%
\lhead{Nome/i Cognome/i}
%

\begin{document}

\begin{exercise}(Martina Apicella)
  Let $\phi(x\,;z)\in L$ and $a\in\U^z$.
  Prove that the following are equivalent
  \begin{itemize}
  \item[1.] the type $\{\phi(x\,;fa)\,:\,f\in\Aut(\U)\}$ is realized (in $\U$)
  \item[2.] there is a formula $\psi(z)\in L$ such that  $\psi(a)$ and $\E x\A z\,[\psi(z)\imp\phi(x\,;z)]$.
  \end{itemize}  
\end{exercise}

\begin{exercise}(Beatrice Calderaro)
  Let $L$ be a countable language containing the language of groups and assume $\U$ is a group.
  Let $p(x)\subseteq L$ define a subgroup of $\U$.
  Prove that there are some formulas $\phi_n(x)\in L$ such that 
  \begin{itemize}
  \item[1.] $p(x)\ \iff\ \{\phi_n(x):n<\omega\}$
  \item[2.] $\phi_{n+1}(x)\wedge\phi_{n+1}(y)\ \imp\ \phi_n(x\cdot y)$.
  \end{itemize}
\end{exercise}

\begin{exercise}(Francesco Dam)
  Let $L$ be a countable language containing the language of groups and assume $\U$ is a group.
  Prove that the following are equivalent
  \begin{itemize}
    \item[1.] every model $M$ is a normal subgroup of $\U$
    \item[2.] some $\omega$-saturated model $M$ is a normal subgroup of $\U$
    \item[3.] $\U$ is a {\sc bfc} group.
  \end{itemize}
  A group $G$ is {\sc bfc} (has boundedly finite conjugacy classes) if for some $n$ the sets $\{g\,a\,g^{-1}:g\in G\}$, as $a$ ranges over $G$, have at most $n$ elements. 
\end{exercise}

\begin{exercise}(Filippo Giugie)
  Let $A\subseteq B$ and $p(x)\subseteq L(A)$. 
  Suppose that $\tp(a/B)$ is isolated over $B$ for every $a\models p(x)$.
  Prove that $p(x)$ is isolated over $A$.
\end{exercise}

\begin{exercise}(Mattia Lo Mele)
   Let $\D\subseteq\U$ be invariant over $A$.
   We say that $\D$ is quasi-definable from below if for every $a\in\D$ and every $\phi(x)\in\tp(a/A)$ there is formula $\psi(x)\in L(A)$ such that $\psi(x)\imp\phi(x)$ and $\psi(\U^x)\subseteq\D$.
   We say that $\D$ is quasi-definable from above if $\neg\D$ is quasi-definable from below.
   Finally $\D$ is quasi-definable if it is both quasi-definable from below and from above.

   Let $L(A)$ be countable. 
   Prove that if $\D$ is quasi-definable then there is a type-definable set $\C$ and a model $M\supseteq A$ such that $M\cap\C=M\cap\D$.
\end{exercise}


\begin{exercise}(Daniele Lanciotti)
Let $p(x)\subseteq L(B)$ and $p_n(x)\subseteq L(A)$, for $n<\omega$, be such that

\hfil$\displaystyle p(x)\imp\bigvee_{i<\omega}p_i(x)$

Prove that 

\hfil$p(x)\wedge\phi(x)\imp p_n(x)$

for some $n<\omega$ and some formula $\phi(x)\in L(A)$ consistent with $p(x)$.
\end{exercise}

\begin{exercise}(Andrea Morello)
  Work in a monster model $\U$.
  Suppose that every formula $\phi(x)\in L$ is equivalent to a some quantifier-free formula $\psi(x)\in L(\U)$.
  Prove that $T$ is $\Delta$-model complete for $\Delta$ the set of formulas of the form $\A y\,\E z\ \theta(x,y,z)$ with $\theta(x,y,z)\in L_{\rm qf}$.
\end{exercise}


\begin{exercise}(Gabriele Maiolino)
  Let $x$ be a finite tuple.
  A set $\B\subseteq\U^x$ is \textit{nowhere dense\/} in the topology induced by $A$ if for every consistent formula $\phi(x)\in L(A)$ there is some consistent $\psi(x)\in L(A)$ such that $\psi(x)\imp\phi(x)$ and $\psi(\B)=\varnothing$.
  A set $\C\subseteq\U^x$ is \textit{meager\/} if it is union of countably many nowhere dense sets.
  
  Assume $L(A)$ is countable.
  Prove that for every countable model $M$ the following are equivalent
  \begin{itemize}
    \item [1.] $\C$ is meager in the topology topology induced by $A$
    \item [2.] there is a model $M$ disjoint of $\C$
  \end{itemize}
\end{exercise}

\end{document}


